%% Merged into Context Engineering Layer subsection of architecture.tex
%% Contains: Tool Result Optimization, Dual-Memory Architecture

\subsubsection{Tool Result Optimization}
\label{sec:tool_results}

Raw tool outputs consume far more tokens than their informational value warrants. A single \texttt{read\_file} may return 2{,}000--3{,}000 tokens of source code, a directory listing may enumerate hundreds of entries, and a test-runner invocation may produce thousands of lines of TAP output. Left unchecked, verbose results dominate the context window within a few iterations, crowding out the user query and system instructions that drive the agent's behavior. Tool result optimization addresses this by transforming raw outputs into compact, semantically preserving representations before they enter the conversation history.

\paragraph{Per-tool-type summarization.}
Each tool result passes through a dedicated summarizer that maps the tool name and raw output to a concise summary (typically 50--200 characters). The summarizer is dispatched by tool name and applies a type-specific compression strategy:

\begin{packeditemize}
\item \textbf{File reads} are replaced with metadata: ``\texttt{$\checkmark$ Read file (142 lines, 4{,}831 chars)}''. The full content remains available through re-reading, but the context carries only the proof that the read occurred and the file's approximate size.

\item \textbf{Search results} report match counts rather than matched lines: ``\texttt{$\checkmark$ Search completed (23 matches found)}''. When a search returns no results, the summary reflects this explicitly so the agent can redirect.

\item \textbf{Directory listings} are collapsed to item counts: ``\texttt{$\checkmark$ Listed directory (47 items)}''. The raw listing, which often contains deeply nested paths, is replaced by a single-line summary.

\item \textbf{Command execution} adapts to output length: short outputs ($\leq$100 characters) are retained verbatim; longer outputs report line counts: ``\texttt{$\checkmark$ Command executed (312 lines of output)}''.

\item \textbf{Errors} are truncated to 200 characters with a classified prefix: ``\texttt{$\times$ Error: FileNotFoundError: ...}''. This ensures the agent receives enough information for error recovery (\Cref{sec:error_recovery}) without consuming excessive context on stack traces.
\end{packeditemize}

\paragraph{Large output offloading.}
For outputs that exceed 8{,}000 characters (approximately 2{,}000 tokens), the summarizer is insufficient: the full output would still dominate the context even after summarization. These outputs are \emph{offloaded} to scratch files before entering the conversation history. The system writes the full output to a session-specific scratch directory (\texttt{\textasciitilde/.opendev/scratch/<session\_id>/}) and replaces the output in the conversation with a 500-character preview plus a reference path: ``\texttt{[Output offloaded: 2{,}341 lines, 48{,}203 chars $\to$ <path>]. Use read\_file to see full output if needed.}'' This creates a natural tiering system: the agent sees enough to understand the content and can \texttt{read\_file} the full output on demand, which is itself subject to the same offloading threshold.

\paragraph{Agent-aware truncation hints.} When output is offloaded to a scratch file, the truncation message includes a tailored recovery hint based on the current agent's capabilities. If the agent has access to subagent delegation (the \texttt{spawn\_subagent} tool), the hint suggests: ``\emph{Delegate to a Code Explorer subagent to process the full output via search and read tools.}'' If the agent lacks subagent capability (e.g., it \emph{is} a subagent), the hint instead suggests: ``\emph{Use the search tool with offset/limit parameters to process the output incrementally.}'' This agent-aware advice prevents the common failure mode where an agent attempts a recovery strategy unavailable in its tool set, for example, a Code Explorer subagent trying to spawn another subagent.

\paragraph{Interaction with compaction.}
Tool result summaries and offloaded outputs serve complementary roles. At ingestion, they provide immediate context savings by replacing the full result in the conversation history. During compaction (\Cref{sec:adaptive_compaction}), the compactor preferentially uses pre-computed summaries when sanitizing messages for LLM-based summarization, avoiding redundant re-processing. This synergy means that even when emergency compaction is triggered, the input to the summarizer LLM is already substantially compressed, improving both the speed and quality of the compaction output.

\paragraph{Design evolution.} Early versions stored full tool output in the conversation history regardless of length. A single long-running test suite could consume 30{,}000 tokens of context in one tool call. The per-tool summarizer reduced this to under 100 tokens in most cases. Adding the 8{,}000-character offloading threshold addressed the remaining outliers (large file reads, verbose command outputs) that exceeded the summarizer's compression ratio, extending typical session length from 15--20 turns (before context overflow) to 30--40 turns without compaction.

\subsubsection{Dual-Memory Architecture for Bounded Thinking}
\label{sec:dual_memory}

The thinking phase (Phase~1 of the ReAct loop, \Cref{sec:react_executor}) requires conversation context for strategic reasoning, but the full conversation history can grow to hundreds of thousands of tokens. Providing the thinking model with unbounded history is infeasible: it would exceed the model's context window and waste budget on stale detail. Providing only recent messages loses strategic context, causing the agent to ``forget'' its overall goals. We resolve this tension through a dual-memory architecture inspired by human cognitive science, which separates compressed long-range context from detailed short-range context.

\paragraph{Episodic memory.}
An LLM-generated summary of the full conversation history captures strategic, long-range context: decisions already made, overall goals, key findings, and important file paths. The summarizer is instructed to preserve actionable identifiers (file paths, function names, variable names, error codes) while omitting verbose tool outputs and redundant exchanges. This summary is regenerated periodically (every 5 new messages, controlled by a \texttt{regenerate\_threshold} parameter) rather than on every turn. Periodic regeneration serves two purposes: it amortizes the cost of the summarization call, and it prevents \emph{summary drift}, a phenomenon where iteratively summarizing a summary causes accumulated distortion. By regenerating from the full history, each episodic memory snapshot is a fresh compression rather than a compression of a compression.

\paragraph{Working memory.}
The last several message pairs (the most recent 6 exchanges by default, controlled by \texttt{exclude\_last\_n}) are reproduced verbatim. These recent messages contain the fine-grained operational details needed for immediate decision-making: exact file contents read in the last few turns, specific error messages, precise line numbers, and the outcome of the most recent tool calls. Summarization would destroy exactly the details that matter most for the next action.

\paragraph{Combined injection.} Before each thinking LLM call, the system constructs the thinking context by concatenating: (1)~the episodic memory summary, providing the ``big picture''; (2)~the working memory messages, providing operational detail; and (3)~the current user query. This structure mirrors the distinction between episodic and working memory in cognitive architecture~\cite{baddeley1992working}: episodic memory stores the gist of past experience for long-range planning, while working memory holds detailed, recently acquired information for immediate use. The thinking token budget remains bounded regardless of conversation length, since the episodic summary has a fixed maximum length (500 characters) and the working memory window is constant.

\paragraph{Design evolution.} Early attempts used pure summarization for the entire history, but critical identifiers (file paths, variable names) were lost, causing the agent to reference non-existent files or misname functions. The opposite extreme, using only recent messages, lost long-range strategic context: the agent would ``forget'' the user's original goal after 10 turns. The hybrid architecture addresses both failure modes. We also discovered that summarizing a previous summary (incremental summarization) accumulates errors over multiple rounds; periodic regeneration from the full history corrects this drift.
